\documentclass[12pt]{article}
\usepackage{amsmath}
\usepackage{graphicx,psfrag,epsf}
\usepackage{enumerate}
\usepackage{natbib}
\usepackage{url}

%user supplied additions
\usepackage{amsfonts}
\usepackage{amsthm}
\usepackage{cleveref}

% DON'T change margins - should be 1 inch all around.
\addtolength{\oddsidemargin}{-.5in}%
\addtolength{\evensidemargin}{-.5in}%
\addtolength{\textwidth}{1in}%
\addtolength{\textheight}{1.3in}%
\addtolength{\topmargin}{-.8in}%


\begin{document}

\def\spacingset#1{\renewcommand{\baselinestretch}%
{#1}\small\normalsize} \spacingset{1}


%%%%%%%%%%%%%%%%%%%%%%%%%%%%%%%%%%%%%%%%%%%%%%%%%%%%%%%%%%%%%%%%%%%%%%%%%%%%%%

\title{\bf Supplementary Material: Sensitivity to Sample Size}
\author{Jackson Loper\\
Department of Statistics, University of Michigan\\
and \\
Jeffrey Regier \\
Department of Statistics, University of Michigan}
\maketitle

\section*{Sensitivity to Sample Size}\label{app:sensitivity}

In this appendix, we present additional simulation results to demonstrate how the bounds vary with different sample sizes.  The main text presents results for $n=100$; here we show analogous results for $n=50$ and $n=200$.

\begin{figure}[h]
\begin{center}
\includegraphics[width=5in]{comparisons/dirichlet_50.png}
  \caption{Comparison of bounds for heterogeneous intervals with $n=50$ variables.}\label{fig:dirichlet50}
\end{center}
\end{figure}

\begin{figure}[h]
\begin{center}
\includegraphics[width=5in]{comparisons/homogen_50.png}
  \caption{Comparison of bounds for homogeneous intervals with $n=50$ variables.}\label{fig:homogen50}
\end{center}
\end{figure}

\begin{figure}[h]
\begin{center}
\includegraphics[width=5in]{comparisons/dirichlet_200.png}
  \caption{Comparison of bounds for heterogeneous intervals with $n=200$ variables.}\label{fig:dirichlet200}
\end{center}
\end{figure}

\begin{figure}[h]
\begin{center}
\includegraphics[width=5in]{comparisons/homogen_200.png}
  \caption{Comparison of bounds for homogeneous intervals with $n=200$ variables.}\label{fig:homogen200}
\end{center}
\end{figure}

As the sample size increases, all bounds become tighter, but the relative advantage of the tight Chernoff bound over the general Hoeffding inequality remains substantial across all values of $n$ considered.

\end{document}
